\chapter{From Euclidean Lattices to Light Cones: The LaMET Revolution}

Parton distributions are light-cone correlations; the lattice computes
Euclidean correlators. For thirty years the bridge between them was
limited to Mellin moments of twist-2 operators. In 2013 Xiangdong Ji
proposed computing the full $x$-dependence instead: boost the hadron,
measure an equal-time spatial correlation, and match it perturbatively
onto the light-cone distribution. This chapter collects the seven papers
that define the framework (LaMET) and its coordinate-space twin
(pseudo-PDFs).

% ----------------------------------------------------------------------
\section{The Founding Proposal}
\papercard{Xiangdong Ji}
{Parton physics on a Euclidean lattice}
{Phys.\ Rev.\ Lett.\ \textbf{110}, 262002 (2013)}
{arXiv:1305.1539 | DOI: 10.1103/PhysRevLett.110.262002}
\whyselected{The single most influential paper in the entire citation
network mined for this selection ($\sim$19/25 of the theory corpus; cited
by essentially every PDF/DA/TMD article of the library). It is itself an
original article of this library (\texttt{Parton Physics on Euclidean
Lattice.pdf}), verified to be the PRL text.}
\physics{Define the \emph{quasi-PDF} through the equal-time, spatially
separated bilocal operator with a Wilson line,
\begin{equation}\label{eq:quasipdf}
  Q_\Gamma(z,P^z) = \langle P|\,\bar\psi(z\hat z)\,\Gamma\,
  W(z\hat z,0)\,\psi(0)\,|P\rangle ,
\end{equation}
which is computable on the lattice exactly like any three-point function.
At large $P^z$, light-cone dominance emerges dynamically: a one-loop
analysis gives the factorization
\begin{equation}\label{eq:l ametmatch}
  Q(x,P^z,\mu) = \int_x^\infty \frac{dy}{|x|}\;
  K\!\left(\frac{y}{x},\frac{\mu}{P^z}\right) q(y,\mu)
  + \mathcal{O}\!\left(\frac{\Lambda_{\rm QCD}^2}{(P^z)^2}\right),
\end{equation}
with $K$ perturbatively computable and universal. Ji checked at one loop
that $K$ reproduces the Altarelli--Parisi kernel in the short-distance
limit, sketched extensions to gluons, GPDs, TMDs and Wigner functions,
and stated the power correction as the price of the method.}
\impact{This is the birth certificate of a subfield. Every chapter after
this one either implements, refines, renormalizes or contests formula
(\ref{eq:l ametmatch}). Within the library, the quasi-PDF numerics (Lin
group), pseudo-PDF numerics (HadStruc), TMD program (LPC) and even the
agentic-workflow papers all trace their lineage to this proposal.}
\cardend

% ----------------------------------------------------------------------
\section{LaMET as an Effective Field Theory}
\papercard{Xiangdong Ji}
{Parton physics from large-momentum effective field theory}
{Sci.\ China Phys.\ Mech.\ Astron.\ \textbf{57}, 1407--1412 (2014)}
{arXiv:1404.6680 | DOI: 10.1007/s11433-014-5450-3}
\whyselected{Second-most-cited theory hub ($\sim$15/99): the formal
statement of Large-Momentum Effective Theory as a systematic expansion,
quoted by every later methodological paper of the library.}
\physics{Boosted states behave as if governed by an effective theory whose
expansion parameter is $\Lambda_{\rm QCD}/P^z$: observables factorize as
\begin{equation}
  f(P^z) = \sum_n C_n(P^z,\mu)\otimes f_n(\mu)\,
  \left(\frac{\Lambda_{\rm QCD}}{P^z}\right)^{\!\omega_n},
\end{equation}
where $C_n$ are universal matching coefficients and $f_n$ are matrix
elements of a small set of quasi-operators. The paper articulates what
makes the class ``universal'': the same coefficients apply across
hadrons and processes, power corrections are dominated by known sources,
and infrared divergences cancel between $Q$ and $q$. It also frames the
strategy that later became standard practice --- extract at moderate
$P^z$, renormalize nonperturbatively, match, then extrapolate
$\Lambda/P^z\to0$ with controlled systematics.}
\impact{Together with Ji~2013 this pair defines the language in which the
library's LP3/HadStruc/ETMC numerical papers write their introductions.
The RMP review below is the mature statement of this program.}
\cardend

% ----------------------------------------------------------------------
\section{The First Matching Coefficient}
\papercard{Xiangdong Ji, Jian-Hui Zhang, Yong Zhao}
{More on large-momentum effective theory approach to parton physics}
{Nucl.\ Phys.\ B \textbf{924}, 366--376 (2017)}
{arXiv:1706.07416}
and
\papercard{Xiong, Ji, Zhang, Zhao}
{One-loop matching for parton distributions: Nonsinglet case}
{Phys.\ Rev.\ D \textbf{90}, 014051 (2014)}
{arXiv:1310.7471 | DOI: 10.1103/PhysRevD.90.014051}
\whyselected{The 2014 letter computed the first one-loop matching kernel
for the nonsinglet quasi-PDF ($\sim$12/99 citations); the 2017 companion
($\sim$5/99) systematized the power-counting argument and defended the
approach against early criticism. Both are indispensable steps from
proposal to tool.}
\physics{Working in dimensional regularization, the one-loop analysis
yields the explicit kernel
\begin{equation}
  \tilde K\!\left(\frac{x}{y}\right) = \delta(x-y)
  + \frac{\alpha_s}{2\pi}\,k\!\left(\frac{x}{y}\right),
\end{equation}
with $k$ containing both the AP splitting function and finite terms that
absorb the difference between equal-time and light-cone kinematics ---
including the characteristic cusp anomalous dimension of the Wilson line
through $\ln(P^z/\mu)$. The NPB~924 paper proves that the only power
divergence of the quasi operator resides in the Wilson-line self-energy
(factorizable, hence removable) and argues that Ioffe-time methods obey
the same large-$P^z$ logic, unifying the two extraction routes.}
\impact{These kernels are what turn lattice matrix elements into PDFs in
the library's quark-PDF papers (Liu et al.\ 2018/2020 helicity and
unpolarized extractions use precisely this nonsinglet kernel).}
\cardend

% ----------------------------------------------------------------------
\section{The Pseudo-PDF Framework}
\papercard{A.\ V.\ Radyushkin}
{Quasi-parton distribution functions, momentum distributions, and
pseudo-parton distribution functions}
{Phys.\ Rev.\ D \textbf{96}, 034025 (2017)}
{arXiv:1705.01488 | DOI: 10.1103/PhysRevD.96.034025}
\whyselected{Defining paper of the coordinate-space route ($\sim$12/99):
every HadStruc article of the library builds on it, and the ETMC and Lin
groups quote it as co-foundation.}
\physics{Radyushkin recast the same physics in Ioffe-time space. With
$\nu = m\sigma$ ($m=2\pi/T$, spacing $\sigma$), define the Ioffe-time
distribution
\begin{equation}
  \mathcal{M}(\nu,z^2) = \langle P|\bar\psi(z)\Gamma W(z,0)\psi(0)|P\rangle,
  \qquad q(\nu)=\int_{-1}^{1}\!dx\,e^{i\nu x}q(x).
\end{equation}
The lattice measures $\mathcal{M}$ at spacelike $z^2<0$ --- the
\emph{pseudo}-ITD. Dividing out the $z$-independent pion form-factor-like
piece defines the \emph{reduced} ITD, free of wave-function
renormalization:
\begin{equation}
  \mathcal{R}(\nu,z^2) \equiv \frac{\mathcal{M}(\nu,z^2)}{\mathcal{M}(0,z^2)}
  \;\xrightarrow[|z^2|\to0]{ }\;
  \int_{-1}^1\!dx\,e^{i\nu x}q(x) + \mathcal{O}(z^2 m^2,\nu^2 m^2 z^2/\nu_0^2),
\end{equation}
a short-distance factorization in position space, equivalent at one loop
to LaMET matching but organized entirely through OPE-type reasoning
descended from Balitsky--Braun.}
\impact{The library's reduced-pseudo-ITD extractions (Fan--Lin pion gluon,
Khan et al., Egerer et al., Salas-Chavira et al.) implement exactly this
ratio; the equivalence with quasi-PDFs established by Izubuchi et al.\
(below) makes the two routes two dialects of one theory.}
\cardend

% ----------------------------------------------------------------------
\section{First Lattice Implementation}
\papercard{K.\ Orginos, A.\ Radyushkin, J.\ Karpie, S.\ Zafeiropoulos}
{Lattice QCD exploration of parton pseudo-distribution functions}
{Phys.\ Rev.\ D \textbf{96}, 094503 (2017)}
{arXiv:1706.05373 | DOI: 10.1103/PhysRevD.96.094503}
\whyselected{The first numerical realization of the pseudo-PDF idea
($\sim$12/99): the reference point that every subsequent HadStruc and
comparison plot in the library quotes.}
\physics{On quenched and $N_f=2$ anisotropic ensembles the paper extracts
the reduced pseudo-ITD from boosted pion and nucleon matrix elements,
demonstrates the expected scaling with $\nu$ and $z$, performs the
one-loop matching into $x$-space, and shows consistency with global-fit
PDFs within uncertainties --- establishing feasibility, error budgeting
(excited states, $z^2$ truncation, momentum reach) and the workflow
used ever since.}
\impact{Together with Radyushkin~2017 this paper founded the second
numerical school of lattice parton physics. Its descendants in this
library include the nucleon/pion/kaon gluon pseudo-ITD series and the
polarized explorations.}
\cardend

% ----------------------------------------------------------------------
\section{The Factorization Theorem}
\papercard{T.\ Izubuchi, X.\ Ji, L.\ Jin, I.\ W.\ Stewart, Y.\ Zhao}
{Factorization theorem relating Euclidean and light-cone parton distributions}
{Phys.\ Rev.\ D \textbf{98}, 056004 (2018)}
{arXiv:1801.03917 | DOI: 10.1103/PhysRevD.98.056004}
\whyselected{The rigorous backbone of both routes ($\sim$14/99): quoted
whenever a library paper needs to justify that its Euclidean observable
actually converges to the light-cone PDF.}
\physics{Using the operator product expansion in position space, the
paper proves that the quasi-PDF obeys the LaMET factorization formula
with calculable coefficients, and that the pseudo-PDF satisfies the
equivalent reduced-ITD factorization; the two formulations are related by
Fourier transformation order by order. Two structural results matter for
practice: the matching coefficient depends on the \emph{parton} momentum
$xP^z$, not on the hadron momentum alone --- later formalized as the
natural scale choice $\mu\sim\sqrt{|x|}P^z$ and exploited by
resummation papers in the library --- and the proof exhibits where the
power corrections live, enabling systematic error estimates.}
\impact{After this theorem, ``quasi vs.\ pseudo'' ceased to be competing
claims and became cross-checks. The library's comparison plots between
HadStruc (pseudo) and LP3 (quasi) results rest on this common ground.}
\cardend

% ----------------------------------------------------------------------
\section{The Mature Synthesis}
\papercard{X.\ Ji, Y.-S.\ Liu, Y.\ Liu, J.-H.\ Zhang, Y.\ Zhao}
{Large-momentum effective theory}
{Rev.\ Mod.\ Phys.\ \textbf{93}, 035005 (2021)}
{arXiv:2004.03543 | DOI: 10.1103/RevModPhys.93.035005}
\whyselected{The authoritative review ($\sim$11/99 and rising): the
canonical citation for the whole framework, written when the field's
foundations had hardened.}
\physics{A comprehensive synthesis: quasi-distributions for PDFs, GPDs,
TMDs and distribution amplitudes; the complete renormalization story
(linear divergences, auxiliary fields, hybrid schemes); power corrections
and renormalons; resummation of longitudinal logarithms; and the status
of numerical results circa 2020, including the library's own LP3 program.
It also clarifies conceptual points that individual letters could only
assert: universality classes of quasi-observables, the role of the
Wilson-line cusp, and why moments-based and $x$-dependent methods are
complementary rather than rival.}
\impact{In this library the review functions as the shared textbook: the
transversity and Boer--Mulders papers cite it as their definitional
reference, and the agentic-workflow papers adopt its pipeline as the
domain logic they automate. Reading it first makes every entry of
Chapters~5--6 easier.}
\cardend
